% ----------------------------------------------------------------
\subsection{L'esempio completo}\label{sec:completo}

Rispetto agli esempi
precedenti la \textsc{map} dello stream completo
(Fig.~\ref{fig:completo}) aggiunge due cose da
osservare. La prima è la testa del grano: il glifo non è una freccia ma la
silhouette della finestra, alta quanto la durata del grano; e la finestra è
a sua volta un inviluppo che attraversa tre stati (Hann, \texttt{rexpodec},
Bartlett), sicché la forma della testa cambia lungo il brano. Le due lenti la
rendono leggibile: sagome affusolate e distinte, colorate per rapporto di
pitch. La seconda è la correlazione fra densità, \texttt{offset\_range} e
durata sulla linea di lettura congelata.

Lo stream è una voce sola. L'ambiente ne prevede $N$, indipendenti per
posizione, griglia, altezza e spazio: è lo scostamento successivo, e
l'oggetto della sezione che segue.

\begin{figure}[H]
    \begin{framedfloat}
    \lstinputlisting[
      linerange={7-21},
      language=yaml,
      basicstyle=\ttfamily\fontsize{7}{8.4}\selectfont
    ]{examples/complete_example/complete_example.yml}
    \vspace{2pt}
    \includegraphics[width=\linewidth]{examples/complete_example/complete_example_map}
    \caption{In alto, lo stream completo (\texttt{complete\_example.yml}).
    Sotto, la \textsc{map} del rendering (assi
    come in Fig.~\ref{fig:c-e}). Il pannello superiore mostra i
    grani nel piano tempo--posizione di lettura, colorati per rapporto di
    pitch; il pannello inferiore gli inviluppi dei parametri principali.}
    \label{fig:completo}
    \end{framedfloat}
\end{figure}

